\documentclass[12pt]{article}
\usepackage[utf8]{inputenc}
\usepackage[T1]{fontenc}
\usepackage{hyperref}
\title{Venda Language Resources for K-12 Curriculum Development}
\author{Compiled by Claude Sonnet 5 (claude-sonnet-5) for Andr\'as Kornai}
\date{}
\begin{document}
\maketitle
\section{Introduction}
Venda (Tshivenda) is an official language of South Africa and a language of instruction. This compilation is based on a single search and is therefore limited, mainly to dictionaries and learner references.
\section{Primary Educational Resources}
\begin{itemize}
\item Oxford First Bilingual Dictionary: Tshivenda and English (over 400 words, for Grades 2--4): \url{https://oxford.co.za/shop/resources-and-home-learning/resource-book/9780195988246-oxford-first-bilingual-dictionary-tshivenda-and-english/}
\end{itemize}
\section{Linguistic and Reference Materials}
\begin{itemize}
\item Bilingual Dictionary English--Tshivenda / Tshivenda--English (about 4000 words, 2021): \url{https://clarkesbooks.co.za/products/bilingual-dictionary-english-tshivenda-tshivenda-english}
\item Sowetan article on a new community Tshivenda-English dictionary (2025): \url{https://www.sowetan.co.za/news/2025-04-09-new-venda-english-dictionary-opens-up-world-of-words/}
\end{itemize}
\section{Religious and Community Resources}
None found in this search.
\section{Digital and Audio Resources}
\begin{itemize}
\item Venda language resources link collection (grammar, pronunciation, flashcards): \url{https://universeofmemory.com/venda-language-resources/}
\end{itemize}
\section{Curriculum Development Context}
Venda is already used in education, so national curriculum documents, CAPS home-language materials and Department of Basic Education textbooks would be the main sources; they were not found in this search and should be sought separately.
\end{document}